\section{Modeling Consensus in Lean~4}


We now apply these foundations to model and verify BFT consensus.

\subsection{System Model}

\begin{definition}[BFT System]
A BFT system consists of:
\begin{itemize}[nosep]
  \item A set of $n$ validators, at most $f$ of which are Byzantine ($f < n/3$).
  \item An asynchronous network: messages may be delayed arbitrarily but are eventually delivered.
  \item A state machine replicated across all honest validators.
\end{itemize}
\end{definition}

\begin{lstlisting}[language=Lean,caption={BFT system model in Lean~4}]
-- Network model
inductive Message where
  | propose : Round -> Block -> Message
  | vote : Round -> BlockHash -> ValidatorId -> Message
  | commit : Round -> Certificate -> Message

-- Validator state
structure ValidatorState where
  round : Nat
  lockedRound : Option Nat
  lockedBlock : Option Block
  decisions : List (Round * Block)

-- System state
structure SystemState (n : Nat) where
  validators : Fin n -> ValidatorState
  network : List Message  -- in-flight messages
  faulty : Finset (Fin n)
  faulty_bound : faulty.card < n / 3
\end{lstlisting}

The \texttt{faulty\_bound} field is a proof obligation: any function that constructs a
\texttt{SystemState} must prove that the number of faulty validators is below the
threshold.

\subsection{Safety: No Two Honest Validators Decide Differently}

\begin{theorem}[Consensus Safety]
If honest validators $i$ and $j$ both decide in round $r$, they decide on the same block.
\end{theorem}

The proof proceeds by contradiction. Suppose $i$ decides block $B$ and $j$ decides block
$B' \neq B$ in round $r$. Both decisions require a quorum certificate---a set of
$\lceil 2n/3 \rceil$ votes. Since there are $n$ validators and at most $f < n/3$ are
Byzantine, any two quorums of size $\lceil 2n/3 \rceil$ must overlap in at least one
honest validator. This honest validator voted for both $B$ and $B'$ in round $r$, which
contradicts the voting rule (an honest validator votes for at most one block per round).

\begin{lstlisting}[language=Lean,caption={Safety proof sketch in Lean~4}]
theorem consensus_safety
    (sys : SystemState n)
    (i j : Fin n)
    (hi : i \notin sys.faulty)
    (hj : j \notin sys.faulty)
    (r : Nat)
    (di : decides sys i r B)
    (dj : decides sys j r B')
    : B = B' := by
  -- Both decisions require quorum certificates
  obtain ⟨qi, hqi⟩ := quorum_of_decision di
  obtain ⟨qj, hqj⟩ := quorum_of_decision dj
  -- Two quorums of size >= 2n/3 overlap
  have overlap := quorum_intersection qi qj sys.faulty_bound
  -- The overlap contains an honest validator
  obtain ⟨v, hv_in_qi, hv_in_qj, hv_honest⟩ := overlap
  -- Honest validators vote for one block per round
  have := honest_single_vote hv_honest r hv_in_qi hv_in_qj
  exact this
\end{lstlisting}

\subsection{The Quorum Intersection Lemma}

This is the heart of BFT safety. We prove it from first principles using Lean~4's
arithmetic.

\begin{lemma}[Quorum Intersection]
If $|Q_1| \geq \lceil 2n/3 \rceil$ and $|Q_2| \geq \lceil 2n/3 \rceil$ and $|F| < n/3$,
then $Q_1 \cap Q_2 \setminus F \neq \emptyset$.
\end{lemma}

\begin{lstlisting}[language=Lean,caption={Quorum intersection proof}]
theorem quorum_intersection
    {n : Nat}
    (Q1 Q2 : Finset (Fin n))
    (F : Finset (Fin n))
    (hQ1 : Q1.card >= 2 * n / 3 + 1)
    (hQ2 : Q2.card >= 2 * n / 3 + 1)
    (hF : F.card < n / 3)
    : (Q1 ∩ Q2 \ F).Nonempty := by
  -- By inclusion-exclusion:
  -- |Q1 cap Q2| >= |Q1| + |Q2| - n
  --             >= 2*(2n/3 + 1) - n
  --             = n/3 + 2
  -- |Q1 cap Q2 \ F| >= n/3 + 2 - |F| > 0
  have h1 := Finset.card_union_le Q1 Q2
  have h2 : (Q1 ∩ Q2).card >= n / 3 + 2 := by omega
  have h3 : (Q1 ∩ Q2 \ F).card >= 2 := by
    calc (Q1 ∩ Q2 \ F).card
        >= (Q1 ∩ Q2).card - F.card := Finset.card_sdiff_le_card ..
      _ >= n / 3 + 2 - F.card := by omega
      _ > 0 := by omega
  exact Finset.card_pos.mp (by omega)
\end{lstlisting}

\subsection{Liveness: Eventual Progress}

Safety is easy relative to liveness. Liveness requires showing that the protocol
eventually makes progress despite Byzantine faults and asynchronous delivery.

\begin{theorem}[Consensus Liveness]
If the network is eventually synchronous (there exists a Global Stabilization Time after
which all messages between honest validators are delivered within a known bound $\Delta$),
then all honest validators eventually decide.
\end{theorem}

The proof uses well-founded induction on the round number. After GST, an honest proposer's
proposal reaches all honest validators within $\Delta$. Since honest validators form a
quorum ($n - f \geq 2n/3 + 1$ when $f < n/3$), they produce a quorum certificate within
$2\Delta$ (one round trip). The key lemma: if a round has an honest proposer and the
network is synchronous, the round succeeds.

\begin{lstlisting}[language=Lean,caption={Liveness argument structure}]
theorem consensus_liveness
    (sys : SystemState n)
    (gst : Nat)  -- Global Stabilization Time
    (sync_after_gst : SynchronousAfter sys gst)
    (fair_proposer : EventuallyHonestProposer sys gst)
    : ∀ v, v ∉ sys.faulty -> EventuallyDecides sys v := by
  intro v hv
  -- After GST, find round r with honest proposer
  obtain ⟨r, hr_after_gst, hr_honest_proposer⟩ := fair_proposer
  -- In synchronous round with honest proposer, all honest
  -- validators receive proposal and vote
  have votes := synchronous_round_votes sync_after_gst hr_after_gst
                  hr_honest_proposer
  -- Honest votes form a quorum
  have quorum := honest_quorum_sufficient sys.faulty_bound votes
  -- Quorum certificate leads to decision
  exact decision_from_quorum quorum hv
\end{lstlisting}

%% ============================================================
